\section{Introduction}
\label{sec:intro}

Cross-lingual search and retrieval-augmented generation depend on multilingual dense retrievers that match English queries to documents in dozens of target languages. Late-interaction retrievers, introduced by ColBERT \citep{khattab2020colbert} and refined by ColBERTv2 \citep{santhanam2022colbertv2}, differ from single-vector encoders in a fundamental way: they score each (query, document) pair at the level of individual token vectors via a MaxSim operation. Every query token finds its best matching document token, and the per-token best matches are summed. Recent multilingual extensions such as ColBERT-XM \citep{louis2024colbertxm} and Jina-ColBERT-v2 \citep{jha2024jinacolbert} extend late interaction to multiple languages, alongside dense single-vector retrievers like LaBSE \citep{feng2022labse}, Contriever \citep{izacard2022contriever}, and multilingual E5 \citep{wang2024me5}.

Token-level scoring places token geometry at the center of retrieval quality. Contextual embeddings from pretrained language models are known to be anisotropic: representations of unrelated tokens concentrate in a narrow cone of the hypersphere, with mean pairwise cosine similarity well above zero \citep{ethayarajh2019contextual,mu2018allbutthetop,gao2019representation}. In single-vector retrievers, this anisotropy is partially masked: the score depends on a sentence-pooled representation, so token-level concentration is averaged out before the comparison. In late-interaction retrievers, every token participates in the score through a hard argmax. If many candidate tokens have similar cosine similarity to a query token, the argmax becomes unstable and the score depends weakly on the identity of the true match. We expect the pathology to be most damaging on low-resource cross-lingual pairs, where multilingual encoders inherit the weakest pretrained signal and the retrieval headroom is largest. \citet{rajaee2022isotropy} document anisotropy in multilingual BERT but stop short of evaluating its consequences for retrieval.

Existing remedies for anisotropy are designed for single-vector encoders and either operate post-hoc or pool away token-level structure. BERT-flow \citep{li2020sentenceembeddings} maps sentence embeddings to an isotropic Gaussian with a learned normalizing flow. WhiteningBERT \citep{huang2021whiteningbert} applies a closed-form whitening transform. SimCSE \citep{gao2021simcse} induces partial isotropy as a side effect of dropout-based contrastive learning. Two recent papers extend isotropy mitigation to ColBERT, both in monolingual settings: \citet{jung2023isotropic} apply post-hoc whitening to ColBERT token vectors on English retrieval benchmarks, and \citet{kim2024hil} propose a hybrid isotropy/anisotropy architecture for English zero-shot transfer. The most relevant training-time regularizer, I-STAR \citep{rudman2024istar}, adds a differentiable isotropy penalty to fine-tuning for classification and concludes that \emph{decreasing} isotropy can help. The question of whether a training-time, token-level regularizer can improve multilingual late-interaction retrieval, and whether the direction of the effect agrees with single-vector or contradicts it, remains open.

This paper presents IsoColBERT, a single-term regularizer added to the InfoNCE training objective of multilingual ColBERT. The regularizer penalizes the mean squared off-diagonal cosine similarity of the valid token vectors in each training batch, with one hyperparameter for its weight and zero cost at inference time. Across three random seeds on FLORES+, IsoColBERT improves MaxSim P@1 by $3.45$ percentage points on English-Swahili (every seed individually positive) and by $1.99$ percentage points on English-Arabic, with the gain extending zero-shot to two further low-resource languages (Nepali $+1.46$ pp, Tamil $+0.88$ pp), while leaving high-resource pairs (EN-ES, EN-FR, EN-DE) unchanged. Token-level geometry shifts as designed: intra-cosine drops and effective rank rises, with the largest shift on Spanish tokens, which collapse hardest under bitext contrastive pressure. Mechanism ablations against two alternative placements rule out localized and matching-level regularization as substitutes for the full-manifold objective.

This paper makes the following contributions.
\begin{enumerate}
    \item We identify token-level anisotropy as a measurable pathology in multilingual late-interaction retrieval. Vanilla ColBERT trained with bitext InfoNCE produces token vectors with intra-cosine of $0.29$ and effective rank of $102$ on Spanish tokens, well below the geometry of English tokens trained under the same objective.
    \item We introduce IsoColBERT, a one-term training-time regularizer that improves token geometry monotonically with its weight and recovers significant retrieval accuracy on low-resource language pairs. The method adds zero inference cost and requires no architectural changes.
    \item In contrast to prior work that attributes multilingual retrieval gaps to data scarcity or supervision, and in contrast to recent isotropy regularization results that find \emph{decreasing} isotropy helps classification fine-tuning, we show that for late-interaction retrieval the dominant correctable cause is geometric and the productive direction is toward higher isotropy across the full token manifold.
\end{enumerate}
